% Einheit 3 von 6 – Einheit 3 (bank/funktionsklassen-und-eigenschaften/e3.jsonl, zusammenbau v0.1)
%% TODO Prüfungswort Einheit 3: nicht in der Marken-Zeile
%% TODO Zeitmarke Einheit 3: Marken-Zeile nicht lesbar
%% TODO Zweigzeile Teil 1: was hier gelernt wird (Satz in Schülersprache aus dem Einheitstitel) – Einheit 3
%% TODO Einheitentitel 3 nicht in der Mappe lesbar
\einheitenkopf[e3]{Einheit 3 von 6 · Einheit 3}
\setcounter{aufgabe}{14}

%% TODO Titel: Ich-kann-Satz fehlt in der Bank (Platzhalter: Kettenname) – Nr. 15
%% TODO Anweisung: ein Satz über der ersten Teilaufgabe fehlt in der Bank – Nr. 15
\begin{aufgabe}{Definitionsbereich, Wertemenge und Schranken}
%% TODO \rechenplatz{n}: Schrittzahl des Grundfalls fehlt in der Bank (nur bei fester Schreibform, 2.3 b)
\begin{teile}
\teil $f(x) = \mathrm{ln}(x - 3)$ – größtmöglicher Definitionsbereich? D = \leerfeld
\teil $f(x) = e^x + 2$ – Wertemenge? W = \leerfeld
\teil Begründe: $f(x) = 5 - 2e^{-x}$ bleibt stets unter $5$.
\teil $f(x) = e^{-x^2}$ – Wertemenge? W = \leerfeld
\teil $f(x) = 2 + \mathrm{sin}(x)$ für $0 \le x < 2\pi$ – Wertemenge? Wie oft wird der Wert $2$ angenommen?
\teil Schar $g_k(x) = \mathrm{ln}(k \cdot x^4 + 1)$ mit $k > 0$. Gib den Definitionsbereich an, zeige, dass alle Graphen durch den Ursprung gehen, und bestimme den exakten Wert von $k$ mit $g_k(1) = 3$ \hfill (Abitur 2017 LK).
\end{teile}
\end{aufgabe}

%% TODO Titel: Ich-kann-Satz fehlt in der Bank (Platzhalter: Kettenname) – Nr. 16
%% TODO Anweisung: ein Satz über der ersten Teilaufgabe fehlt in der Bank – Nr. 16
\begin{aufgabe}{Definitionsbereich, Wertemenge und Schranken – Fehler finden, Begründen, Anwendung}
\begin{teile}
\teil Ich finde den Fehler: $f(x) = 4 - e^{-x}$ hat die Wertemenge $W = ]-\infty; 4]$.
\teil Begründe: $\mathrm{ln}(x)$ ist für $x \le 0$ nicht definiert.
\teil Ein Kühlschrank kühlt nach dem Einschalten: $T(t) = 4 + 16e^{-0{,}5t}$ ($T$ in °C, $t$ in Stunden). Lebensmittel sollen erst bei höchstens $5$ °C hinein. Ab wann geht das? Werden jemals $4$ °C erreicht?
\end{teile}
\end{aufgabe}
